\subsection{Alternatif Format dan Struktur PGHD}\label{analisis-alternatif-format-pghd}

Berdasarkan dasar teori pada Subbab~\ref{standar-format-pghd} dan analisis masalah (subbab~\ref{analisis-permasalahan}), terdapat beberapa model dan pertimbangan yang bisa digunakan sebagai format data PGHD untuk DecMed.
Tabel~\ref{tab:perbandingan-format-pghd} merangkum perbandingan antarmodel dan skema format data.
Parameter yang perlu menjadi pertimbangan untuk memilih standar format PGHD di antaranya,

\begin{enumerate}
  \item \textbf{Ukuran payload}: besar payload JSON untuk satu atau sekumpulan pengukuran mempertimbangkan \textit{gas fee} IOTA, penyimpanan IPFS dan efisiensi pengiriman.
  \item \textbf{Kesesuaian dengan Android}: tingkat kemudahan integrasi format dengan ekosistem Android (Health Connect API, Kotlin), termasuk apakah tersedia \textit{library} atau \textit{mapping} bawaan.
  \item \textbf{Preservasi informasi klinis}: kemampuan format menyimpan nilai numerik, satuan, dan konteks temporal yang diperlukan untuk interpretasi klinis dan penelitian lanjutan.
  \item \textbf{Kompleksitas implementasi}: mempertimbangkan batasan implementasi di Android (subbab~\ref{batasan-masalah}), beban pengembangan untuk membuat \textit{serializer/deserializer} dan validasi skema di sisi klien Android.
  \item \textbf{Kemampuan \textit{batching}}: kesesuaian format untuk mengelompokkan banyak pengukuran ke dalam satu payload, sesuai dengan strategi batching periodik yang dipilih pada Subbab~\ref{alternatif-arsitektur-pipeline-pghd}.
\end{enumerate}


\begin{table}[!ht]
  \centering
  \caption{Perbandingan format data PGHD untuk pipeline DecMed}\label{tab:perbandingan-format-pghd}
  \footnotesize  % smaller than \small
  \setlength{\tabcolsep}{3.5pt}  % reduced spacing
  \begin{tabular}{lp{2.2cm}p{2.2cm}p{1.8cm}p{2.2cm}p{1.8cm}}
    \toprule
    \textbf{Format} &
    \textbf{Ukuran Payload} &
    \textbf{Kesesuaian Android} &
    \textbf{Preservasi Klinis} &
    \textbf{Kompleksitas} &
    \textbf{Batching} \\
    \midrule
    Apple HealthKit
      & Moderat
      & Tidak ada\textsuperscript{a}
      & Tinggi
      & Tinggi
      & Tidak didukung \\[0.2cm]
    Google Health Connect
      & Moderat
      & Sangat tinggi\textsuperscript{b}
      & Tinggi
      & Rendah
      & Terbatas\textsuperscript{c} \\[0.2cm]
    Open mHealth / IEEE~1752
      & Rendah-sedang
      & Sedang\textsuperscript{d}
      & Tinggi
      & Sedang
      & Didukung (data series) \\[0.2cm]
    HL7 FHIR (Kanta PHR)
      & Tinggi\textsuperscript{e}
      & Sedang\textsuperscript{f}
      & Sangat tinggi
      & Tinggi
      & Tidak dioptimalkan \\[0.2cm]
    \bottomrule
    \multicolumn{6}{l}{\footnotesize\textsuperscript{a}Tidak tersedia di Android.} \\
    \multicolumn{6}{l}{\footnotesize\textsuperscript{b}Native SDK Kotlin, dokumentasi resmi, dukungan WorkManager.} \\
    \multicolumn{6}{l}{\footnotesize\textsuperscript{c}Batch melalui \texttt{insertRecords\(\)}, tetapi sebagai format penyimpanan akhir tidak dirancang untuk IPFS.} \\
    \multicolumn{6}{l}{\footnotesize\textsuperscript{d}Tidak ada SDK Android, perlu implementasi manual.} \\
    \multicolumn{6}{l}{\footnotesize\textsuperscript{e}Verbose karena elemen \texttt{coding}, \texttt{system}, \texttt{display}, dan metadata FHIR.} \\
    \multicolumn{6}{l}{\footnotesize\textsuperscript{f}Tersedia HAPI FHIR library, tetapi ukuran \textit{dependency} besar.} \\
  \end{tabular}
\end{table}

Berdasarkan perbandingan-perbandingan skema yang ada, Apple HealthKit tidak relevan untuk penelitian ini karena implementasi dilakukan di Android, sehingga tidak ada jalur pengambilan data dari HealthKit.

Google Health Connect cocok dengan implementasi berbasis Android karena merupakan API resmi Google untuk akses data kesehatan di Android.
Google Health Connect juga memiliki SDK Kotlin dan panduan penulisan data \parencite{android_hc_write,android_hc_dataformat}.
Namun, format internalnya berbasis kelas Kotlin (\texttt{Record} dan subkelasnya) dan tidak secara langsung menghasilkan payload JSON yang dapat disimpan di IPFS \parencite{android_hc_dataformat}.
Penggunaan Health Connect tetap bisa digunakan sebagai \textit{sumber data}, tetapi memerlukan konversi jika ingin digunakan sebagai format penyimpanan akhir di DecMed.

Open mHealth dan IEEE~1752 memiliki kelebihan kesederhanaan dan standarisasi format data \parencite{openmhealth_1752_blog,ieee_1752_1_2021_standard}
Namun, tidak tersedianya SDK Android berarti implementasi serialisasi, pengolahan data dan validasi skema harus dilakukan secara manual.

HL7 FHIR di Kanta PHR memberikan interoperabilitas dan didukung standardisasi organisasi internasional, tetapi ukuran payload yang besar akibat metadata FHIR yang \textit{verbose} menjadi kendala signifikan untuk penyimpanan di IPFS dalam skala batch data time-series yang besar.
Panduan HL7 PGHD juga menjelaskan bahwa banyak item PGHD konsumen belum sepenuhnya tercakup oleh kode internasional, dan saat ini masih menggunakan kode sementara.\parencite{hl7_2025_pghd_ig}

Berdasarkan analisis di atas, akan digunakan format data hibrid menggunakan Google Health Connect sebagai sumber data di sisi Android, dan merancang format JSON \textit{custom} yang terinspirasi prinsip Open mHealth sebagai format payload yang disimpan di IPFS\@.

% Alasan utama pemilihan ini adalah:
% \begin{enumerate}
%   \item Health Connect adalah antarmuka standar Android untuk data
%     kesehatan; menggunakannya sebagai sumber data memastikan akses
%     ke seluruh tipe data yang kompatibel dengan ekosistem Android secara
%     seragam.\parencite{android_hc_write}
%   \item Format JSON kustom yang terinspirasi Open mHealth mempertahankan
%     prinsip desain utama standar tersebut (nilai + unit + metadata
%     minimum dalam satu objek, \textit{timestamp} ISO~8601) tanpa overhead
%     metadata FHIR yang tidak diperlukan untuk operasional
%     pipeline.\parencite{openmhealth_1752_blog,ieee_1752_1_2021_standard}
%   \item Desain kustom memungkinkan struktur batching yang dioptimalkan
%     untuk pipeline DecMed, yaitu satu payload JSON yang memuat banyak
%     pengukuran dalam satu jendela waktu, sehingga jumlah transaksi
%     IPFS dan IOTA tetap terkendali.
%   \item Di masa mendatang, payload JSON ini dapat dipetakan ke resource
%     FHIR \texttt{Observation} dengan menambahkan lapisan transformasi,
%     karena field utama (tipe, nilai, unit, waktu) selaras secara
%     semantik dengan model FHIR.\parencite{hl7_2025_pghd_ig}
% \end{enumerate}

% \begin{table}[htbp]
%   \centering
%   \caption{Perbandingan format data PGHD untuk pipeline DecMed}
%   \label{tab:perbandingan-format-pghd}
%   \resizebox{\textwidth}{!}{%
%   \begin{tabular}{lcccccc}
%     \toprule
%     \textbf{Format} &
%     \textbf{Interopera-bilitas} &
%     \textbf{Overhead Payload} &
%     \textbf{Kesesuaian Android} &
%     \textbf{Preservasi Klinis} &
%     \textbf{Kompleksitas Impl.} &
%     \textbf{Batching} \\
%     \midrule
%     Apple HealthKit
%       & Rendah\textsuperscript{a}
%       & Moderat
%       & Tidak ada\textsuperscript{b}
%       & Tinggi
%       & Tinggi
%       & Tidak didukung \\[0.2cm]
%     Google Health Connect
%       & Sedang\textsuperscript{c}
%       & Moderat
%       & Sangat tinggi\textsuperscript{d}
%       & Tinggi
%       & Rendah
%       & Terbatas\textsuperscript{e} \\[0.2cm]
%     Open mHealth / IEEE~1752
%       & Tinggi\textsuperscript{f}
%       & Rendah–sedang
%       & Sedang\textsuperscript{g}
%       & Tinggi
%       & Sedang
%       & Didukung (data series) \\[0.2cm]
%     HL7 FHIR (Kanta PHR)
%       & Sangat tinggi
%       & Tinggi\textsuperscript{h}
%       & Sedang\textsuperscript{i}
%       & Sangat tinggi
%       & Tinggi
%       & Tidak dioptimalkan \\[0.2cm]
%     JSON Kustom (DecMed)
%       & Sedang\textsuperscript{j}
%       & Rendah
%       & Tinggi\textsuperscript{k}
%       & Tinggi
%       & Rendah
%       & Dioptimalkan \\
%     \bottomrule
%     \multicolumn{7}{l}{\footnotesize\textsuperscript{a}Proprietary iOS.
%     \textsuperscript{b}Tidak tersedia di Android.
%     \textsuperscript{c}Hanya di ekosistem Android, tidak ada format pertukaran lintas sistem.} \\
%     \multicolumn{7}{l}{\footnotesize\textsuperscript{d}Native SDK Kotlin, dokumentasi resmi, dukungan WorkManager.
%     \textsuperscript{e}Batch melalui \texttt{insertRecords()}, tetapi sebagai format penyimpanan akhir tidak dirancang untuk IPFS.} \\
%     \multicolumn{7}{l}{\footnotesize\textsuperscript{f}Standar IEEE terbuka, tersedia library OMH-to-FHIR.
%     \textsuperscript{g}Tidak ada SDK Android resmi, perlu implementasi manual.} \\
%     \multicolumn{7}{l}{\footnotesize\textsuperscript{h}Verbose karena elemen \texttt{coding}, \texttt{system}, \texttt{display}, dan metadata FHIR.
%     \textsuperscript{i}Tersedia HAPI FHIR library, tetapi overhead dependensi besar.} \\
%     \multicolumn{7}{l}{\footnotesize\textsuperscript{j}Dapat dipetakan ke FHIR atau OMH bila diperlukan.
%     \textsuperscript{k}Serialisasi via \texttt{kotlinx.serialization} atau Gson tanpa dependensi besar.} \\
%   \end{tabular}}
% \end{table}